\section{Batching Optimization}
\label{sec:batching}

Batching decides how strings are grouped before they reach the model. This chapter tests each part of the batching logic
separately: batch layout and sorting, beam size, and parallel execution with an output length cap.

\subsection{Benchmark Method}
A benchmark tool runs inside the production container, with the same model, the same two pinned cores and the same
machine type, and times only model inference (tokenize, translate, detokenize). Its baseline step calls the production
translation function itself, and a sanity check confirmed that the benchmark engine reproduces production output exactly
(100 of 100 strings).

\subsection{Batch Layout and Sorting}
Table~\ref{tab:layout} applies the layout changes to the \code{.resx} workload (1,071 translations). A run-wide plan with
real token counts and 1,024-token batches, sorted by length, cut padding from 30.1\% to 4.3\%. Smaller and larger token
budgets were also tested.

\begin{table}[H]
\centering
\caption[Batch layout and beam size]{Batch layout and beam size on the \code{.resx} workload (run
\code{ec2-20261016t103748z}). ``Identical'' compares the output with the baseline's output.}
\label{tab:layout}
\small\setlength{\tabcolsep}{4.5pt}
\begin{tabular}{@{}l l c S[table-format=3.2] S[table-format=1.2] S[table-format=2.1] c c@{}}
\toprule
\textbf{Step} & \textbf{Change} & \textbf{Beam} & {\textbf{Time (s)}} & {\textbf{Speed-up}} & {\textbf{Padding (\%)}} & \textbf{Peak mem.} & \textbf{Identical} \\
\midrule
A & Original pipeline & 5 & 434.84 & 1.00 & 30.1 & 1.8 GB & -- \\
D-256 & Sorted, 256-token batches & 5 & 451.10 & 0.96 & 1.2 & 1.8 GB & 100\% \\
D-512 & Sorted, 512-token batches & 5 & 432.62 & 1.01 & 2.4 & 2.2 GB & 100\% \\
D-1024 & Sorted, 1,024-token batches & 5 & 425.78 & 1.02 & 4.3 & 3.5 GB & 100\% \\
\midrule
F & D-1024 + beam 2 & 2 & 188.08 & 2.31 & 4.3 & 3.5 GB & 77\% \\
G & D-1024 + beam 1 & 1 & 103.58 & 4.20 & 4.3 & 3.5 GB & 59\% \\
\bottomrule
\end{tabular}
\end{table}

\begin{figure}[H]
\centering
\begin{tikzpicture}
\begin{axis}[reportbars, width=0.92\linewidth, height=5.6cm, bar width=17pt,
  symbolic x coords={A,D-256,D-512,D-1024,F,G}, xtick={A,D-256,D-512,D-1024,F,G},
  ylabel={Inference time (s)}, ymin=0, ymax=520,
  nodes near coords, nodes near coords align={vertical},
  xticklabels={{A\\beam 5},{D-256\\beam 5},{D-512\\beam 5},{D-1024\\beam 5},{F\\beam 2},{G\\beam 1}},
  xticklabel style={align=center}]
\addplot[fill=cprev!80, draw=cprev] coordinates {(A,434.84) (D-256,451.10) (D-512,432.62) (D-1024,425.78)};
\addplot[fill=cnew!80, draw=cnew] coordinates {(F,188.08) (G,103.58)};
\end{axis}
\end{tikzpicture}
\caption[Inference time per batching step]{Inference time for each batching step. Layout changes (red) moved time by at
most 4\%; beam size (blue) halved it or more.}
\label{fig:layout}
\end{figure}

\paragraph{Finding.} Layout changes never altered a translation, but moved time by only $-4\%$ to $+2\%$, within the
2\% run-to-run noise measured on identical work. Padding fell from 30\% to 1--4\% without a matching speed-up:
\textbf{padding is not the bottleneck}. Peak memory, however, roughly doubled with 1,024-token batches, and a 4,096-token
attempt did not complete on the 8~GB machine, so batches above about 1,024 tokens were ruled out. Telling CTranslate2
explicitly to use two threads per batch also changed nothing.

\subsection{Beam Size}
The beam size is the number of candidate translations the model keeps while decoding. Steps F and G show that it, not
the layout, dominates the cost: beam~2 was 2.31\x{} and beam~1 4.20\x{} faster than the baseline. Because a smaller
beam changes the output, quality was checked on two datasets.

\paragraph{FLORES-200.} On 40 human-translated sentences per language (Table~\ref{tab:flores}), beam~2 lost 0.87 chrF and
beam~1 lost 1.40 chrF overall.

\begin{table}[H]
\centering
\caption[chrF on FLORES-200 by beam size]{chrF on FLORES-200 (first 40 dev sentences per language) by beam size.}
\label{tab:flores}
\begin{tabular}{@{}c S[table-format=2.2] *{7}{S[table-format=2.2]}@{}}
\toprule
\textbf{Beam} & {\textbf{All 7}} & {de-DE} & {es-ES} & {fr-FR} & {hi-IN} & {ja-JP} & {ru-RU} & {zh-CN} \\
\midrule
5 & 56.69 & 62.71 & 55.54 & 68.75 & 57.65 & 33.83 & 58.35 & 28.58 \\
2 & 55.82 & 61.93 & 55.63 & 67.05 & 56.98 & 34.77 & 56.71 & 26.32 \\
1 & 55.29 & 60.08 & 55.73 & 67.76 & 55.03 & 33.39 & 57.33 & 25.07 \\
\bottomrule
\end{tabular}
\end{table}

\paragraph{KDE4.} Forty sentences per language is a small sample, so the comparison was repeated on a fixed random sample
of 1,000 KDE4 strings in all twelve languages (12,000 translations per beam), using run-wide, length-sorted 1,024-token
batches (Tables~\ref{tab:kde4beam} and~\ref{tab:kde4beamlang}).

\begin{table}[H]
\centering
\caption[Beam size on the KDE4 sample]{Beam size on 1,000 KDE4 strings $\times$ 12 languages (run \code{ec2-20261016t115802z}).}
\label{tab:kde4beam}
\begin{tabular}{@{}c S[table-format=4.0] S[table-format=5.0] S[table-format=1.2] S[table-format=2.2] c@{}}
\toprule
\textbf{Beam} & {\textbf{Time (s)}} & {\textbf{Translations/hour}} & {\textbf{Speed-up}} & {\textbf{chrF (all 12)}} & \textbf{Identical to beam 5} \\
\midrule
5 & 3169 & 13631 & 1.00 & 64.41 & -- \\
\textbf{2} & \bfseries 1442 & \bfseries 29949 & \bfseries 2.20 & \bfseries 63.96 & 84\% \\
1 & 915 & 47197 & 3.46 & 62.26 & 69\% \\
\bottomrule
\end{tabular}
\end{table}

\begin{table}[H]
\centering
\caption[chrF per language by beam size]{chrF per language on the KDE4 sample by beam size.}
\label{tab:kde4beamlang}
\small
\begin{tabular}{@{}c *{12}{S[table-format=2.2]}@{}}
\toprule
\textbf{Beam} & {hi} & {ar} & {bn} & {de} & {es} & {fr} & {ja} & {mr} & {pa} & {ta} & {te} & {zh} \\
\midrule
5 & 66.90 & 57.15 & 64.87 & 65.06 & 72.01 & 70.23 & 53.43 & 59.34 & 63.35 & 64.70 & 66.07 & 39.68 \\
2 & 66.61 & 55.74 & 64.11 & 64.39 & 71.69 & 69.48 & 52.52 & 59.08 & 63.27 & 65.07 & 65.86 & 39.30 \\
1 & 65.60 & 53.77 & 63.88 & 62.21 & 68.49 & 68.16 & 46.51 & 58.76 & 62.68 & 63.70 & 64.26 & 33.72 \\
\bottomrule
\end{tabular}
\end{table}

\begin{figure}[H]
\centering
\begin{tikzpicture}
\begin{axis}[width=0.78\linewidth, height=6.2cm, axis lines*=left, ymajorgrids, xmajorgrids, grid style={dashed,gray!35},
  xlabel={Speed-up over beam 5}, ylabel={chrF}, xmin=0.6, xmax=4.0, ymin=54, ymax=66,
  tick label style={font=\footnotesize}, label style={font=\small},
  legend style={font=\footnotesize, draw=none, legend columns=2, at={(0.5,-0.24)}, anchor=north}]
\addplot[mark=*, thick, color=cnew, nodes near coords, point meta=explicit symbolic,
  every node near coord/.append style={font=\scriptsize, anchor=south west}]
  coordinates {(1.00,64.41)[beam 5] (2.20,63.96)[beam 2] (3.46,62.26)[beam 1]};
\addlegendentry{KDE4 (12 languages)}
\addplot[mark=square*, thick, color=cprev, nodes near coords, point meta=explicit symbolic,
  every node near coord/.append style={font=\scriptsize, anchor=north west}]
  coordinates {(1.00,56.69)[beam 5] (2.15,55.82)[beam 2] (3.65,55.29)[beam 1]};
\addlegendentry{FLORES-200 (7 languages)}
\end{axis}
\end{tikzpicture}
\caption[Speed against quality by beam size]{Speed against quality for beam sizes 5, 2 and 1. Beam 2 keeps almost all
of the quality at more than twice the speed.}
\label{fig:beam}
\end{figure}

\paragraph{Finding.} On real localization text, beam~2 is \textbf{2.2\x{} faster for a 0.45-point chrF loss}, and every
language stays within 1.4 points of beam~5 (Tamil slightly improves). Beam~1 costs more on KDE4 than on FLORES, with large
drops for Japanese ($-6.9$) and Chinese ($-6.0$), which makes it unsuitable for these languages. Beam~2 was adopted for
production.

\subsection{Parallel Batches and an Output Length Cap}
During the KDE4 runs the machine's CPU was only about 59\% busy: CTranslate2 processed one batch at a time and part of the
two cores sat idle. Three steps were tested on 500 KDE4 strings in 12 languages (6,000 translations, beam~2):

\begin{itemize}
  \item \textbf{P1:} load the model with two engine workers (\code{inter\_threads=2}, \code{intra\_threads=1}) and queue all
        batches at once, so two batches are translated in parallel, one per core.
  \item \textbf{P2:} P1 with smaller 512-token batches.
  \item \textbf{P3:} P1 plus a per-batch output cap of twice the longest input plus 10 tokens. Because batches are sorted by
        length, the cap fits every string in the batch and only stops runaway outputs.
\end{itemize}

\begin{table}[H]
\centering
\caption[Parallel batches and output length cap]{Parallel batches and output length cap (run \code{ec2-20261016t142418z}, 6,000 translations, beam 2).}
\label{tab:parallel}
\begin{tabular}{@{}l l S[table-format=3.2] S[table-format=1.2] S[table-format=2.1] c c@{}}
\toprule
\textbf{Step} & \textbf{Change} & {\textbf{Time (s)}} & {\textbf{Speed-up}} & {\textbf{CPU busy (\%)}} & \textbf{Identical} & \textbf{Peak mem.} \\
\midrule
P0 & One batch at a time & 617.52 & 1.00 & 59.3 & -- & 2.6 GB \\
P1 & Two batches in parallel & 514.97 & 1.20 & 99.6 & 100\% & 4.5 GB \\
P2 & P1 with 512-token batches & 539.68 & 1.14 & 99.3 & 99.9\% & 4.5 GB \\
\textbf{P3} & \textbf{P1 + output length cap} & \bfseries 488.14 & \bfseries 1.27 & 99.6 & \textbf{100\%} & 4.5 GB \\
\bottomrule
\end{tabular}
\end{table}

\paragraph{Finding.} Parallel batches raised CPU use from 59\% to 99.6\% and made translation 1.20\x{} faster with
identical output. The gain is below 2\x{} because two cores sharing the work each run slower than one core alone, and
it costs a second set of working buffers (2.6 to 4.5~GB, still safe on 8~GB). Smaller batches were slower, because more
batches mean more per-batch overhead. The length cap added a further 6\% (1.27\x{} in total) without changing a single
translation.

\subsection{Adopted Configuration}
\begin{table}[H]
\centering
\caption[Adopted batching configuration]{Batching configuration before and after this chapter.}
\label{tab:adopted}
\begin{tabularx}{\linewidth}{@{}l L L@{}}
\toprule
\textbf{Setting} & \textbf{Before} & \textbf{After} \\
\midrule
Order of strings & file order & sorted by length within each file, in chunks of 1,000 strings \\
Batch size & 52 strings or $\sim$1,024 \emph{estimated} tokens & up to 1,024 \emph{real} tokens \\
Execution & one batch at a time & all batches of a chunk queued, two engine workers \\
Output length & unlimited & 2 $\times$ longest input + 10 tokens per batch \\
Beam size & 5 & 2 \\
\bottomrule
\end{tabularx}
\end{table}

\keyresult{\textbf{RQ2.} Padding and batch layout are not the bottleneck: cutting padding from 30\% to 4\% gained 3\%.
The real levers are the beam size (2.2\x{} for $-0.45$ chrF on KDE4) and keeping both cores busy (1.27\x{} with identical
output). Together these are about 2.8\x{} faster in the benchmark.}
